% ECONOMICS_PROBLEM: {"id": "D74-356", "title": "Conflict, fragility and poverty traps", "jel_code": "D74", "source_id": "OP-0039"}
% FIRST_POSTED: 2026-10-09
% ATTEMPT_STATUS: unresolved
% ENTRY_KIND: research_attempt
\documentclass[11pt]{article}
\usepackage[T1]{fontenc}
\usepackage[utf8]{inputenc}
\usepackage[margin=25mm]{geometry}
\usepackage{amsmath,amssymb,amsthm,xurl,hyperref}
\newtheorem{proposition}{Proposition}
\newtheorem{lemma}{Lemma}
\newtheorem{corollary}{Corollary}
\newcommand{\E}{\mathbb E}
\newcommand{\R}{\mathbb R}
\newcommand{\Prb}{\mathbb P}
\newcommand{\Var}{\operatorname{Var}}
\DeclareMathOperator*{\argmax}{arg\,max}
\DeclareMathOperator*{\argmin}{arg\,min}
\setlength{\parindent}{0pt}
\setlength{\parskip}{6pt}
\newcommand{\Cov}{\operatorname{Cov}}
\title{Conflict, fragility and poverty traps: a research attempt}
\author{GPT-6 (Codex)}
\date{9 October 2026}
\begin{document}
\maketitle

\section*{Target and outcome}
\textbf{Status: unresolved.} The desired event requires a first entry by horizon $H$ followed by $J$ uninterrupted periods in a good region. It is neither eventual entry nor stationary prosperity. This attempt gives an exact finite-state control construction, a sharp distinction between those targets, and a model-error bound. It does not identify causal reconstruction transitions from actual conflict data.

Blattman and Miguel's primary review abstract \cite{review} highlights the policy and identification gaps surrounding conflict and recovery. The dynamic program below is an editorial mathematical specialization, not a reconstruction theorem attributed to that review.

\section*{A finite model with explicit memory}
Let $S$ be a finite set, $G\subseteq S$ the predeclared target region, and $U(s)$ finite nonempty policy sets. Each action already incorporates a complete feasible security, transfer and services package. Specify a controlled transition matrix $P_u(s,s')$. The state includes remaining funds and political constraints whenever these affect feasible actions. Assume observed state and Markov sufficiency; hidden political conditions would instead require a belief-state formulation.

A direct recursion on physical state alone is inadequate. After first entry into $G$, an early exit makes the specified event fail permanently, even if the economy later recovers. A dynamic program that resets its clock after that exit would solve a different event. Augment the state by one of: not yet entered; entered with $j$ required transitions remaining; permanently failed; or completed durability. This finite memory makes the event a terminal success indicator.

Define $D_0(s)=\mathbf1\{s\in G\}$ and, for $j\geq1$,
\[
 D_j(s)=\mathbf1\{s\in G\}\max_{u\in U(s)}
                         \sum_{s'}P_u(s,s')D_{j-1}(s').
\]
This is the maximal chance of remaining in $G$ for the next $j$ transitions conditional on having just entered it. The indicator prevents an exit followed by reentry from being counted as success. Time-dependent transitions can be handled by adding calendar time to $D$.

For physical states $s\notin G$ with no previous entry, put $A_H(s)=0$ and recurse backwards for $t=H-1,\ldots,0$:
\[
 A_t(s)=\max_{u\in U(s)}\sum_{s'}P_u(s,s')
 \left[\mathbf1\{s'\in G\}D_J(s')+
       \mathbf1\{s'\notin G\}A_{t+1}(s')\right].
\]
If the initial state is already in $G$, the answer is $D_J(s_0)$ because first entry is at time zero. Otherwise it is $A_0(s_0)$. When $J=0$, first entry itself is enough; when $H=0$ and $s_0\notin G$, the answer is zero.

\begin{proposition}[Exact finite-state solution]
The displayed recursion equals the supremum of the catalogue's durability probability over all admissible history-dependent rules. An optimal deterministic rule depending on the augmented state exists.
\end{proposition}
\begin{proof}
Conditional on a current augmented state and chosen action, the next-state probabilities are $P_u$. The terminal payoff is zero or one. Backward induction over the at most $H+J$ transitions gives the conditional success value. At each node the maximum over a finite nonempty action set is attained. Randomization averages action values and cannot improve their maximum. Past history has no predictive or feasibility content beyond the augmented state by assumption. Combining nodewise maximizing actions therefore attains the history-dependent supremum.
\end{proof}
With dense transitions and at most $A$ actions, computation takes $O((H+J)|S|^2A)$ arithmetic operations, plus whatever cost is needed to enumerate budget states. This is useful only if the finite state representation itself is manageable.

\section*{A two-state diagnostic}
Let the physical states be bad $B$ and good $G$, with a fixed policy producing
\[
 P(B,G)=a,\qquad P(G,B)=b,\qquad0<a,b<1.
\]
Starting at $B$, first entry occurs at date $k$ with probability $(1-a)^{k-1}a$. Conditional on entry, survival for $J$ transitions has probability $(1-b)^J$. Thus
\[
 p(a,b;H,J)=[1-(1-a)^H](1-b)^J.
\]
The stationary good-state mass is instead $a/(a+b)$, obtained by equating flow $(1-v)a=vb$. Both stationary probabilities are positive, so this chain has no absorbing poverty class. A long spell of poverty in such a model is persistence, not proof of an absorbing trap.

For $H=1,J=5$, policy A with $(a,b)=(0.8,0.4)$ has success probability $0.8(0.6)^5=0.062208$, while policy B with $(0.2,0.01)$ has probability $0.2(0.99)^5\approx0.190198$. A larger short-run recovery probability can therefore be inferior for durable recovery. The numerical values are synthetic, not empirical conflict estimates.

The formula also separates useful margins:
\[
 \partial_a p=H(1-a)^{H-1}(1-b)^J,\qquad
 \partial_b p=-J[1-(1-a)^H](1-b)^{J-1}.
\]
A reconstruction package can improve entry while worsening durability, for example by rapidly deploying resources without sustainable institutional support. These derivatives state a conditional tradeoff; that narrative is not an identified mechanism without independent evidence.

\section*{Identification, welfare and robust error}
Sequential randomization, consistency and positivity identify transition laws for supported packages when conditioning uses sufficient pretreatment history. If assignment depends on unobserved violence risk, simply tabulating observed transitions estimates selection as well as policy effects. Spillovers can be included by taking $S$ to represent a group of linked regions; ignoring cross-border displacement generally changes the intervention being identified.

Suppose each estimated transition row differs from the true row by at most $\delta$ in total variation, uniformly over states and feasible actions. Couple true and fitted paths under the same history-dependent rule while their histories agree. At each step the chance of a new mismatch is at most $\delta$. Hence for every rule,
\[
 |p_P(u)-p_{\widehat P}(u)|\leq \min\{1,(H+J)\delta\}.
\]
The same bound holds for optimal values by taking suprema. A fitted-model optimal rule loses at most $2(H+J)\delta$ in true success probability, capped at one, by comparing its fitted and true values and those of a true optimizer. A confidence statement needs a simultaneous data-based guarantee for $\delta$.

Welfare requires a separate reward function including income, violence harms and real program costs with fixed weights. Maximizing success probability alone may spend excessive resources on crossing an arbitrary threshold. One can impose an expected budget constraint using finite-horizon occupation measures on the augmented state, or scalarize success and welfare with a stated weight. Expected constraints can require policy randomization even though the unconstrained success problem does not. Cash transfers within the social boundary cancel except for distributional weights and financing distortions.

\section*{Remaining work}
The recursion is an exact solution after transitions, feasible actions and sufficient states are supplied. Those primitives are the empirical substance of the original agenda. Establishing sequencing effects, regional spillovers and the existence of trapping classes remains unresolved. A practical study should report first-entry, uninterrupted-survival and welfare outcomes separately rather than labeling all three as escape from a poverty trap.

\section{Completion Estimate}
48\%

This is a post hoc, heuristic estimate for the full catalogue target.
An augmented-state dynamic program exactly optimizes finite-horizon durable recovery, and coupling bounds quantify sensitivity to transition error. Causal reconstruction kernels, sequencing and spillover effects, and welfare comparisons remain to be established from actual post-conflict data.

\begin{thebibliography}{9}
\bibitem{review} C. Blattman and E. Miguel (2010), Civil War, \emph{Journal of Economic Literature} 48(1), 3--57. Primary publisher abstract and metadata inspected. \url{https://www.aeaweb.org/articles?id=10.1257/jel.48.1.3}.
\end{thebibliography}

\end{document}
